%!TEX root = ../main.tex
\chapter{Introduction \`a R pour la Statistique}
\label{ch:r_intro}

\begin{intuition}
Si Python est un couteau suisse g\'en\'eraliste, \textbf{R}\index{R (langage)} est un
instrument de pr\'ecision con\c{c}u sp\'ecifiquement pour la statistique. Cr\'e\'e par des
statisticiens pour des statisticiens, R excelle dans l'analyse de donn\'ees, les tests
d'hypoth\`eses et la visualisation. De nombreux chercheurs en biologie, \'ecologie et
sciences sociales l'utilisent quotidiennement.
\end{intuition}

% ----------------------------------------------------------------
\section{Pourquoi R ?}

\begin{definition}[Le langage R]
\textbf{R} est un langage de programmation libre et gratuit, sp\'ecialis\'e dans le calcul
statistique et la repr\'esentation graphique. Il dispose d'un \'ecosyst\`eme de plus de
20\,000 paquets sur le d\'ep\^ot CRAN\index{CRAN}.
\end{definition}

\begin{center}
\begin{tikzpicture}[
  langbox/.style={draw, rounded corners, minimum width=3.5cm, minimum height=2.8cm,
                  align=center, font=\small},
  feat/.style={font=\footnotesize, align=left},
]
  \node[langbox, fill=blue!60, text=white] (python) at (0,0)
    {\textbf{\Large Python}};
  \node[feat, below=0.2cm of python.south, anchor=north] {
    $\bullet$ G\'en\'eraliste\\
    $\bullet$ Machine learning\\
    $\bullet$ Automatisation\\
    $\bullet$ NumPy, Pandas, SciPy
  };

  \node[langbox, fill=green!50!black, text=white] (r) at (7,0)
    {\textbf{\Large R}};
  \node[feat, below=0.2cm of r.south, anchor=north] {
    $\bullet$ Statistique native\\
    $\bullet$ Tests d'hypoth\`eses\\
    $\bullet$ Visualisation (ggplot2)\\
    $\bullet$ Bioconductor, tidyverse
  };

  \draw[<->, thick, >=stealth] (python.east) -- (r.west)
    node[midway, above, font=\bfseries] {Compl\'ementaires};
\end{tikzpicture}
\end{center}

% ----------------------------------------------------------------
\section{Variables et vecteurs}
\index{vecteur (R)}\index{c() (R)}

En R, l'op\'erateur d'affectation usuel est \ri{<-} (fl\`eche). La structure de base
est le \textbf{vecteur}, cr\'e\'e avec la fonction \ri{c()}.

\begin{codeblock}[title=\textbf{R}]
\begin{minted}{r}
# Affectation
x <- 42
nom <- "physique"

# Vecteur num\'erique
mesures <- c(12.3, 14.1, 11.8, 13.5, 12.9)
print(mesures)

# Op\'erations vectoris\'ees
mesures_cm <- mesures * 100
print(mesures_cm)
\end{minted}
\end{codeblock}

\begin{codeoutput}
\begin{verbatim}
[1] 12.3 14.1 11.8 13.5 12.9
[1] 1230 1410 1180 1350 1290
\end{verbatim}
\end{codeoutput}

\begin{attention}
En R, l'indexation commence \`a \textbf{1} (et non 0 comme en Python).
Ainsi, \ri{mesures[1]} renvoie le \emph{premier} \'el\'ement, soit 12.3.
\end{attention}

\begin{codeblock}[title=\textbf{R}]
\begin{minted}{r}
# Indexation (commence `a 1 !)
print(mesures[1])    # Premier \'el\'ement
print(mesures[2:4])  # \'El\'ements 2 `a 4 (inclus !)

# S\'equences et r\'ep\'etitions
indices <- 1:10
print(indices)
repetitions <- rep(0, 5)
print(repetitions)
\end{minted}
\end{codeblock}

\begin{codeoutput}
\begin{verbatim}
[1] 12.3
[1] 14.1 11.8 13.5
 [1]  1  2  3  4  5  6  7  8  9 10
[1] 0 0 0 0 0
\end{verbatim}
\end{codeoutput}

% ----------------------------------------------------------------
\section{Data frames en R}
\index{data.frame (R)}

\begin{codeblock}[title=\textbf{R}]
\begin{minted}{r}
# Cr\'eer un data frame
experience <- data.frame(
  sujet = c("A", "B", "C", "D", "E"),
  dose  = c(10, 20, 30, 40, 50),
  effet = c(2.1, 4.5, 5.8, 7.2, 8.9)
)
print(experience)
str(experience)
\end{minted}
\end{codeblock}

\begin{codeoutput}
\begin{verbatim}
  sujet dose effet
1     A   10   2.1
2     B   20   4.5
3     C   30   5.8
4     D   40   7.2
5     E   50   8.9
'data.frame':	5 obs. of  3 variables:
 $ sujet: chr  "A" "B" "C" "D" ...
 $ dose : num  10 20 30 40 50
 $ effet: num  2.1 4.5 5.8 7.2 8.9
\end{verbatim}
\end{codeoutput}

% ----------------------------------------------------------------
\section{Fonctions statistiques de base}
\index{mean (R)}\index{sd (R)}\index{var (R)}\index{median (R)}\index{summary (R)}

R int\`egre nativement toutes les fonctions statistiques essentielles.

\begin{codeblock}[title=\textbf{R}]
\begin{minted}{r}
notes <- c(12, 15, 8, 14, 16, 11, 13, 9, 17, 10)

cat("Moyenne     :", mean(notes), "\n")
cat("M\'ediane    :", median(notes), "\n")
cat("\'Ecart-type :", sd(notes), "\n")
cat("Variance    :", var(notes), "\n")
cat("Min / Max   :", min(notes), "/", max(notes), "\n")

# R\'esum\'e complet
summary(notes)
\end{minted}
\end{codeblock}

\begin{codeoutput}
\begin{verbatim}
Moyenne     : 12.5
Mediane     : 12.5
Ecart-type  : 3.027650
Variance    : 9.166667
Min / Max   : 8 / 17
   Min. 1st Qu.  Median    Mean 3rd Qu.    Max.
   8.00   10.25   12.50   12.50   14.75   17.00
\end{verbatim}
\end{codeoutput}

% ----------------------------------------------------------------
\section{Tests statistiques}
\index{t.test (R)}\index{cor.test (R)}

\begin{definition}[Test de Student]
Le \textbf{test t de Student}\index{test t} compare les moyennes de deux groupes
pour d\'eterminer si la diff\'erence observ\'ee est statistiquement significative
(p-valeur $< 0.05$).
\end{definition}

\begin{codeblock}[title=\textbf{R}]
\begin{minted}{r}
# Deux groupes : traitement vs contr\^ole
traitement <- c(23.1, 25.4, 22.8, 26.1, 24.5, 27.0)
controle   <- c(19.2, 20.1, 18.5, 21.3, 19.8, 20.5)

resultat <- t.test(traitement, controle)
print(resultat)
\end{minted}
\end{codeblock}

\begin{codeoutput}
\begin{verbatim}
	Welch Two Sample t-test

data:  traitement and controle
t = 5.1893, df = 9.5182, p-value = 0.0004507
alternative hypothesis: true difference in means is not equal to 0
95 percent confidence interval:
 2.633574 6.566426
sample estimates:
mean of x mean of y
 24.81667  20.21667
\end{verbatim}
\end{codeoutput}

\begin{codeblock}[title=\textbf{R}]
\begin{minted}{r}
# Test de corr\'elation
x <- c(1, 2, 3, 4, 5, 6, 7, 8)
y <- c(2.1, 4.3, 5.8, 8.2, 9.7, 12.1, 13.8, 16.2)
cor.test(x, y)
\end{minted}
\end{codeblock}

\begin{codeoutput}
\begin{verbatim}
	Pearson's product-moment correlation

data:  x and y
t = 42.085, df = 6, p-value = 4.386e-08
95 percent confidence interval:
 0.9971346 0.9999227
sample estimates:
      cor
0.9984112
\end{verbatim}
\end{codeoutput}

% ----------------------------------------------------------------
\section{R\'egression lin\'eaire avec \ri{lm()}}
\index{lm() (R)}\index{r\'egression lin\'eaire}

\begin{codeblock}[title=\textbf{R}]
\begin{minted}{r}
# R\'egression lin\'eaire : effet en fonction de la dose
modele <- lm(effet ~ dose, data = experience)
summary(modele)
\end{minted}
\end{codeblock}

\begin{codeoutput}
\begin{verbatim}
Coefficients:
            Estimate Std. Error t value Pr(>|t|)
(Intercept)  0.14000    0.34351   0.408  0.71062
dose         0.17200    0.01033  16.649  0.00048 ***
---
Residual standard error: 0.3742 on 3 degrees of freedom
Multiple R-squared:  0.9893,	Adjusted R-squared:  0.9857
\end{verbatim}
\end{codeoutput}

Le coefficient $R^2 = 0.989$ indique un excellent ajustement lin\'eaire :
l'effet augmente de 0.172 unit\'es par unit\'e de dose.

% ----------------------------------------------------------------
\section{Graphiques de base en R}
\index{plot (R)}\index{hist (R)}\index{boxplot (R)}

\begin{codeblock}[title=\textbf{R}]
\begin{minted}{r}
# Nuage de points avec droite de r\'egression
plot(experience$dose, experience$effet,
     xlab = "Dose (mg)", ylab = "Effet",
     main = "Relation dose-effet", pch = 19, col = "blue")
abline(modele, col = "red", lwd = 2)

# Histogramme
hist(rnorm(1000), breaks = 30, col = "lightblue",
     main = "Distribution normale simul\'ee",
     xlab = "Valeur")

# Bo\^ite `a moustaches
boxplot(traitement, controle,
        names = c("Traitement", "Contr\^ole"),
        col = c("coral", "lightgreen"),
        main = "Comparaison des groupes")
\end{minted}
\end{codeblock}

% ----------------------------------------------------------------
\section{Comparaison Python vs R}

\begin{exemple}[M\^eme analyse dans les deux langages]

\begin{codeblock}[title=\textbf{Python}]
\begin{minted}{python}
import numpy as np
from scipy import stats

data = [12, 15, 8, 14, 16, 11, 13, 9, 17, 10]
print(f"Moyenne : {np.mean(data):.2f}")
print(f"Ecart-type : {np.std(data, ddof=1):.2f}")

# Test t
groupe_a = [23.1, 25.4, 22.8, 26.1, 24.5, 27.0]
groupe_b = [19.2, 20.1, 18.5, 21.3, 19.8, 20.5]
t_stat, p_val = stats.ttest_ind(groupe_a, groupe_b)
print(f"p-value : {p_val:.6f}")
\end{minted}
\end{codeblock}

\begin{codeblock}[title=\textbf{R}]
\begin{minted}{r}
data <- c(12, 15, 8, 14, 16, 11, 13, 9, 17, 10)
cat("Moyenne :", mean(data), "\n")
cat("Ecart-type :", sd(data), "\n")

# Test t
groupe_a <- c(23.1, 25.4, 22.8, 26.1, 24.5, 27.0)
groupe_b <- c(19.2, 20.1, 18.5, 21.3, 19.8, 20.5)
resultat <- t.test(groupe_a, groupe_b)
cat("p-value :", resultat$p.value, "\n")
\end{minted}
\end{codeblock}
\end{exemple}

\begin{remarque}
Notez que \ri{sd()} en R calcule l'\'ecart-type \emph{corrig\'e} ($n-1$) par d\'efaut,
comme \py{np.std(data, ddof=1)} en Python. La fonction \py{np.std()} sans
\py{ddof=1} divise par $n$, ce qui donne un r\'esultat diff\'erent.
\end{remarque}

% ----------------------------------------------------------------
\begin{errmark}{Erreurs fr\'equentes en R}
\begin{enumerate}
  \item \textbf{Indexation \`a partir de 1}\index{indexation R} --- En R, \ri{x[1]} est le
    \emph{premier} \'el\'ement. En Python, \py{x[1]} est le \emph{deuxi\`eme}. Cette diff\'erence
    est source de nombreuses erreurs quand on passe d'un langage \`a l'autre.
  \item \textbf{\ri{<-} vs \ri{=}} --- En R, on privil\'egie \ri{<-} pour l'affectation.
    Le signe \ri{=} fonctionne aussi dans la plupart des cas, mais \ri{<-} est la
    convention standard. Attention \`a ne pas \'ecrire \ri{x < -5} (comparaison !) au
    lieu de \ri{x <- 5} (affectation).
  \item \textbf{Facteurs inattendus} --- R convertit parfois automatiquement les cha\^ines
    en \texttt{factor}. Utilisez \ri{stringsAsFactors = FALSE} dans
    \ri{data.frame()} si n\'ecessaire.
  \item \textbf{Oublier \ri{na.rm = TRUE}} --- Les fonctions comme \ri{mean()} retournent
    \ri{NA} si le vecteur contient des valeurs manquantes. Ajoutez \ri{na.rm = TRUE}.
\end{enumerate}
\end{errmark}

% ----------------------------------------------------------------
\begin{projet}{Analyse statistique comparative en R}
R\'ealisez une \'etude compl\`ete en R :
\begin{enumerate}
  \item Cr\'eez un data frame avec deux groupes exp\'erimentaux (15 mesures chacun),
    simul\'es avec \ri{rnorm(15, mean=..., sd=...)}.
  \item Calculez les statistiques descriptives de chaque groupe (\ri{mean}, \ri{sd},
    \ri{median}).
  \item R\'ealisez un test t de Student pour comparer les deux groupes.
  \item Tracez des bo\^ites \`a moustaches c\^ote \`a c\^ote.
  \item Ajoutez une troisi\`eme variable et effectuez une r\'egression lin\'eaire avec \ri{lm()}.
  \item Interpr\'etez les r\'esultats : la diff\'erence est-elle significative ?
    Le mod\`ele lin\'eaire est-il pertinent ($R^2$) ?
\end{enumerate}
\end{projet}

% ----------------------------------------------------------------
\section{Exercices}

\begin{exercice}[$\star$]
Cr\'eez un vecteur contenant les temp\'eratures moyennes mensuelles d'une ville
(12 valeurs). Calculez la moyenne, la m\'ediane et l'\'ecart-type. Quel mois est le
plus chaud ? Utilisez \ri{which.max()}.
\end{exercice}

\begin{exercice}[$\star$]
Cr\'eez un data frame avec les colonnes \ri{element}, \ri{symbole} et
\ri{masse\_atomique} pour 5 \'el\'ements chimiques. Affichez le r\'esum\'e avec
\ri{summary()} et acc\'edez \`a la colonne \ri{masse\_atomique} avec \ri{\$}.
\end{exercice}

\begin{exercice}[$\star\star$]
G\'en\'erez 100 valeurs al\'eatoires suivant une loi normale ($\mu = 50$, $\sigma = 10$)
avec \ri{rnorm()}. Tracez un histogramme et superposez la courbe th\'eorique avec
\ri{curve(dnorm(x, 50, 10), add = TRUE)}.
\end{exercice}

\begin{exercice}[$\star\star$]
Simulez une exp\'erience : g\'en\'erez deux \'echantillons de taille 30, l'un de moyenne
100 et l'autre de moyenne 105 (m\^eme \'ecart-type 15). Effectuez un test t.
R\'ep\'etez 1000 fois et comptez le pourcentage de fois o\`u $p < 0.05$
(c'est la \textbf{puissance}\index{puissance statistique} du test).
\end{exercice}

\begin{exercice}[$\star\star\star$]
Chargez le jeu de donn\'ees int\'egr\'e \ri{iris} en R. R\'ealisez une analyse compl\`ete :
statistiques descriptives par esp\`ece, test ANOVA avec \ri{aov()} pour comparer les
longueurs de s\'epales entre esp\`eces, et r\'egression lin\'eaire entre longueur et
largeur de p\'etales. Produisez 3 graphiques pertinents.
\end{exercice}

% ----------------------------------------------------------------
\begin{formules}{Chapitre 9 --- R : l'essentiel}
\begin{tabular}{@{}p{5cm} p{9cm}@{}}
\toprule
\textbf{Concept} & \textbf{Syntaxe R} \\
\midrule
Affectation & \ri{x <- valeur} \\
Vecteur & \ri{c(1, 2, 3)} \\
S\'equence & \ri{1:10} ou \ri{seq(0, 1, by=0.1)} \\
Data frame & \ri{data.frame(col1=..., col2=...)} \\
Moyenne & \ri{mean(x)} \\
\'Ecart-type & \ri{sd(x)} \\
R\'esum\'e & \ri{summary(x)} \\
Test t & \ri{t.test(groupe1, groupe2)} \\
Corr\'elation & \ri{cor.test(x, y)} \\
R\'egression & \ri{lm(y \~{} x, data=df)} \\
Graphique & \ri{plot(x, y)}, \ri{hist(x)}, \ri{boxplot(x)} \\
\bottomrule
\end{tabular}
\end{formules}
